\documentclass[11pt]{article}

\usepackage[utf8]{inputenc}
\usepackage[T1]{fontenc}
\usepackage{lmodern}
\usepackage[margin=1in]{geometry}
\usepackage{microtype}
\usepackage{enumitem}
\usepackage{hyperref}
\usepackage{xurl}

\hypersetup{
  colorlinks=true,
  linkcolor=blue,
  citecolor=blue,
  urlcolor=blue,
  pdftitle={Formal Verification of a Runtime Security Boundary for Model Context Protocol Tool Calls},
  pdfauthor={Paolo Vella}
}

\title{Formal Verification of a Runtime Security Boundary for Model Context Protocol Tool Calls}
\author{Paolo Vella\\Vellaveto\\\texttt{security@vellaveto.online}}
\date{March 9, 2026}

\begin{document}

\maketitle

\begin{abstract}
Runtime mediation for AI agents is security-critical because the mediator
decides whether side-effecting tool invocations may proceed. A mediator that
fails open ceases to be a meaningful security control. This paper presents the
formal verification program for Vellaveto, a runtime security boundary for
Model Context Protocol (MCP) tool calls and adjacent agent actions. The
verification stack combines TLA+ model checking~\cite{lamport94,lamport02},
Alloy bounded relational modeling~\cite{jackson12}, Lean~4 and Coq theorem
proving~\cite{lean4,coq}, Verus deductive verification over production Rust
kernels~\cite{verus}, and Kani bounded model checking over production Rust
code~\cite{kani}. The objective is not a monolithic proof of the entire system,
but a layered proof architecture for the narrow waist of the security design:
fail-closed verdict computation, delegation attenuation, approval safety,
audit-chain integrity, and normalization kernels that could otherwise create
bypass channels. The repository further exposes a checked trusted-assumption
inventory, an explicit refinement map from abstract policy models to traced
Rust evaluation, and pinned reproducibility entrypoints~\cite{vellavetoFormalInventory,vellavetoTCB,vellavetoRefinementMap,vellavetoFormalDocker}. The resulting
claim is intentionally narrow: the security-critical pure logic at the runtime
boundary is covered by a complementary proof mesh, and the remaining trust
assumptions are explicit.
\end{abstract}

\noindent\textbf{Keywords:} formal verification, MCP, agent security, policy
enforcement, Rust, Verus, Kani, TLA+, audit integrity, capability delegation

\section{Introduction}

Model Context Protocol servers and agent frameworks increase the set of actions
large-language-model-based systems can take on the world~\cite{mcpSpec}. Those
actions include filesystem access, network calls, command execution, resource
reads, and other operations with persistent or externally visible consequences.
The runtime mediator between an agent and the tool or service it invokes is
therefore a natural enforcement point for security policy.

That boundary is a difficult verification target. It is not enough to show that
the policy language is expressive, or that common attacks are covered by a test
suite. The crucial question is whether the implementation can make the wrong
decision under edge conditions:

\begin{itemize}[leftmargin=*]
\item Can an unmatched request ever produce \texttt{Allow}?
\item Can a delegated capability become broader than the parent that issued it?
\item Can approval tokens be replayed across sessions or across different actions?
\item Can normalization or decoding corner cases create a bypass around a deny rule?
\item Can audit verification silently accept malformed state and thereby
undermine later integrity claims?
\end{itemize}

This paper argues that such a boundary benefits from a layered proof strategy.
Abstract state-machine models are necessary, but insufficient on their own. The
implementation must also be targeted directly, especially the pure kernels that
compute verdicts, validate structure, and project sensitive context.

The work described here is implemented in the \texttt{formal/} subtree of the
Vellaveto repository~\cite{vellavetoFormalInventory} and targets the runtime
security logic in the Vellaveto workspace. The verified claims are narrower
than ``the whole system is proved,'' but they cover the logic whose failure
would most directly invalidate the system's security contract.

The paper makes three main claims. First, the repository contains a
multi-method verification architecture in which each proof technology is used
for the class of property it models well. Second, the strongest proof layer is
implementation-facing: Verus and Kani target production Rust kernels rather
than a separate executable specification. Third, the artifact states its proof
boundary explicitly through a refinement map and a checked trusted-assumption
inventory, rather than leaving that boundary implicit.

\section{Problem Statement}

An MCP security boundary is only meaningful if it fails closed and preserves
monotonicity across derived security state. In practice, that means at least
the following properties:

\begin{enumerate}[leftmargin=*]
\item Missing policy, missing context, malformed derived state, or explicit deny
conditions must not produce \texttt{Allow}.
\item Capability delegation must only attenuate privilege, never expand it.
\item Approval gates must be tied to the intended action and session and must
not be reusable after successful consumption.
\item Audit verification must reject malformed chains, malformed Merkle proofs,
inconsistent rotation manifests, and malformed structured decision metadata.
\item Parser, normalization, and decoding helpers must not create alternate
interpretations that widen access beyond what policy intended.
\end{enumerate}

These are implementation-facing safety obligations. They cannot be discharged
by benchmark numbers, API documentation, or integration tests alone.

\section{Scope}

The proof scope is the runtime security boundary implemented by the Vellaveto
workspace. At a high level, requests arrive through one of several transport
surfaces, are normalized into actions, are evaluated against policy and
security-relevant context, and yield a verdict plus audit metadata.

The proof-relevant subsystems include:

\begin{itemize}[leftmargin=*]
\item policy evaluation and rule override behavior in \texttt{vellaveto-engine},
\item capability delegation, deputy validation, and context projection in
\texttt{vellaveto-mcp},
\item approval scope binding and approval consumption gates,
\item audit append and verification logic in \texttt{vellaveto-audit},
\item normalization kernels such as path normalization, IDNA handling, and DLP
buffer arithmetic, and
\item selected state machines such as task lifecycle, credential-vault
behavior, and cascading-failure transitions.
\end{itemize}

The proof inventory is maintained in the repository's formal index, while the
trusted boundary and non-goals are documented separately~\cite{vellavetoFormalInventory,vellavetoTCB,vellavetoFormalScope}.

\section{Contributions}

This work makes five concrete contributions.

\subsection{A Multi-Method Verification Architecture}

The repository combines TLA+, Alloy, Lean~4, Coq, Verus, and Kani rather than
depending on a single proof technology. Each tool is assigned the class of
claim it models well: state-machine properties, relational attenuation,
abstract semantics, deductive Rust-kernel proofs, or bounded model checking of
concrete code.

\subsection{Implementation-Facing Deductive Proofs over Production Rust Kernels}

The strongest component of the artifact is the Verus suite, which targets
production Rust kernels rather than a separate reference implementation. This
includes verdict computation, constraint evaluation, approval gates, capability
checks, audit verification, Merkle structure, path normalization, DLP
arithmetic, and ACIS envelope validation~\cite{verus,vellavetoFormalInventory}.

\subsection{Bridge Harnesses between Abstract and Concrete Proof Layers}

Kani harnesses are used not only for standalone properties, but also to bridge
preconditions around deductive proofs. For example, bounded model checking is
used to establish ordering properties that support proofs of priority-sensitive
verdict behavior~\cite{kani,vellavetoFormalInventory}.

\subsection{An Explicit Refinement and Trust-Boundary Story}

The repository does not leave correspondence between abstract models and
implementation logic implicit. Instead, it includes a refinement map and a
checked trusted-assumption inventory~\cite{vellavetoRefinementMap,vellavetoTCB}.
This makes the proof boundary inspectable.

\subsection{Reproducibility Infrastructure}

The repository exposes top-level reproduction entrypoints through
\texttt{make formal} and \texttt{make verify}, and ships a pinned Docker image
for replaying the formal environment~\cite{vellavetoFormalDocker}.

\section{Verification Architecture}

\subsection{TLA+}

TLA+ is used for operational state-machine and protocol properties, including
fail-closed policy-engine behavior, ABAC forbid-override semantics, task
lifecycle safety, cascading-failure control flow, capability-delegation
invariants, credential-vault transitions, and audit-chain integrity~\cite{lamport94,lamport02,vellavetoFormalInventory}.

This layer is best suited to properties about allowed transitions and global
invariants rather than line-by-line implementation behavior.

\subsection{Alloy}

Alloy is used for compact bounded relational models, especially around
capability attenuation and ABAC combining logic~\cite{jackson12}. These models
are useful where structure matters more than stepwise control flow.

\subsection{Lean 4 and Coq}

Lean~4 and Coq encode semantic properties over abstract models, including
fail-closed evaluation, determinism, path normalization, ABAC
forbid-override, and delegation attenuation~\cite{lean4,coq}. These proofs
establish clean semantic baselines, but do not by themselves eliminate the
refinement gap to the Rust implementation.

\subsection{Verus}

Verus is the implementation-oriented core of the stack. It proves properties
over extracted or isolated Rust kernels from the production codebase. Examples
include:

\begin{itemize}[leftmargin=*]
\item core verdict computation and rule overrides,
\item constraint-evaluation fail-closed control flow,
\item approval binding and single-use consumption gates,
\item approval-ID validation on transport and server paths,
\item capability matching, subset, grant-coverage, and identity-chain guards,
\item deputy and principal handoff logic,
\item audit append and audit verification guards,
\item Merkle append, fold, path, and rotation-manifest structure,
\item cross-call DLP arithmetic,
\item path normalization, and
\item ACIS envelope validation.
\end{itemize}

This is the layer that most directly supports the claim that the runtime
boundary's security-critical pure logic has proof coverage on actual Rust
code~\cite{verus,vellavetoFormalInventory}.

\subsection{Kani}

Kani complements Verus by exploring bounded concrete execution spaces in actual
Rust code. It is well suited for finite-state and edge-condition properties,
including IP classification, cache safety, grant subset behavior, constraint
edge cases, Unicode and IDNA normalization, RwLock poisoning fail-closed
behavior, sanitizer inversion properties, and cascading-failure finite-state
machine transitions~\cite{kani,vellavetoFormalInventory}.

\section{Representative Verified Properties}

The repository contains many individual proof obligations. For exposition, they
are grouped here by security meaning rather than by file count.

\subsection{Fail-Closed Verdict Computation}

The central safety claim is that unmatched or malformed evaluation state must
not yield \texttt{Allow}. This is supported at multiple levels:

\begin{itemize}[leftmargin=*]
\item TLA+ models fail-closed defaults and priority ordering.
\item Lean~4 and Coq prove abstract fail-closed semantics.
\item Verus proves the core Rust verdict kernel for all inputs.
\item Kani checks bounded concrete harnesses, including error and edge cases.
\end{itemize}

This property is the cornerstone of the runtime boundary. If it fails, the rest
of the security design loses force.

\subsection{Rule Override and Deny Dominance}

Security boundaries often fail not because the base policy is wrong, but
because an override path is mishandled. The proofs therefore isolate path and
network block behavior and ensure that rule-level deny conditions dominate the
final verdict when triggered.

\subsection{Delegation Attenuation}

Capability and identity delegation are first-class escalation surfaces. The
proof stack covers them through abstract attenuation models and concrete Rust
guards for matching, subset checks, grant coverage, identity-chain continuity,
and depth restrictions.

\subsection{Approval Safety}

Approval paths are verified as gates, not merely as workflow artifacts. The
implementation-facing proofs target session binding, action-fingerprint
binding, single-use consumption, and the validation of approval identifiers
carried through transport and HTTP-facing paths.

\subsection{Audit Integrity}

Tamper evidence is structurally verified through audit append and recovery
counters, hash-chain verification guards, Merkle proof structure, rotation
manifest linkage, and structured envelope validation. The verified claim is
structural integrity, not proof of underlying cryptographic primitives.

\subsection{Safe Normalization and Decoding}

Several recent attack classes against agent tooling depend on alternate parsing
or normalization paths. This repository therefore treats helpers such as path
normalization, DLP buffer arithmetic, homoglyph collapse, IDNA fail-closed
handling, and lock-poisoning handlers as proof-relevant kernels.

\section{Refinement Boundary}

This work does not claim a complete machine-checked forward simulation from the
TLA+ models to the full Rust implementation. Instead, it documents the
refinement boundary explicitly in the repository's refinement map~\cite{vellavetoRefinementMap}.

That artifact defines:

\begin{itemize}[leftmargin=*]
\item an abstraction function from traced Rust evaluation to abstract TLA state,
\item per-transition simulation obligations,
\item executable witness tests for selected high-risk safety transitions, and
\item explicit non-claims, including unfinished obligations such as full
stuttering-refinement coverage.
\end{itemize}

This is a deliberate methodological choice. The artifact narrows the gap rather
than obscuring it.

\section{Trusted Assumptions}

The verified claims in this paper are conditional on a named trust boundary.
The relevant assumptions include:

\begin{itemize}[leftmargin=*]
\item compiler and proof-tool correctness,
\item cryptographic primitive correctness,
\item operating-system and filesystem semantics,
\item DNS and network infrastructure behavior, and
\item the assumption that the runtime boundary is actually invoked for the
actions being controlled.
\end{itemize}

The repository treats this as an auditable part of the artifact through the
trusted computing base documentation and the checked assumption inventory under
\texttt{formal/tools/}~\cite{vellavetoTCB}.

\section{Reproducibility}

The artifact can be replayed using the repository entrypoints:

\begin{verbatim}
make formal
make verify
\end{verbatim}

Tool-specific targets:

\begin{verbatim}
make formal-tla
make formal-alloy
make formal-lean
make formal-coq
make formal-kani
make formal-verus
make formal-trusted-assumptions
\end{verbatim}

Pinned Docker environment:

\begin{verbatim}
docker build -t vellaveto-formal formal/
docker run --rm -v "$(pwd):/workspace" vellaveto-formal
\end{verbatim}

The Docker image pins the major formal toolchains used by the artifact,
including Rust, Verus, Kani, TLA+, and Lean tooling~\cite{vellavetoFormalDocker}.

\section{Related Work}

This artifact sits at the intersection of formal methods and runtime security
mediation.

At the methods level, it builds on well-established verification ecosystems:
TLA+ for state-machine model checking~\cite{lamport94,lamport02}, Alloy for
bounded relational reasoning~\cite{jackson12}, Coq and Lean~4 for
machine-checked semantics~\cite{coq,lean4}, and Verus and Kani for
implementation-oriented verification of Rust code~\cite{verus,kani}.

At the application level, the novelty is not a new logic, but the application
of a mixed proof stack to the runtime mediation layer of MCP-style tool calls.
The repository's proof strategy differs from work that verifies a standalone
algorithm in isolation: the target here is a privilege-bearing security
boundary.

\section{Threats to Validity}

Several limitations affect the strength of the overall claim.

\subsection{No End-to-End Proof}

The repository does not contain a full end-to-end mechanical proof of the
entire deployed system. Glue code, network behavior, and external dependencies
remain outside the strongest proof layer.

\subsection{Bounded Tools Remain Bounded}

TLA+ and Kani results are exhaustive only within their configured bounds or
symbolic encodings. They are strong evidence, but they are not substitutes for
unbounded deductive proofs in every domain.

\subsection{Abstract-Model Correspondence Remains Partial}

Lean~4 and Coq proofs target abstract semantics. They strengthen the semantic
story, but do not remove the implementation-refinement problem on their own.

\subsection{Cryptography and Environment Are Assumed}

Cryptographic primitive correctness, operating-system behavior, filesystem
semantics, and DNS or network infrastructure remain trusted assumptions rather
than proved properties.

\section{Conclusion}

This paper presents a formal verification program for a runtime security
boundary rather than for a single algorithm in isolation. The central design
principle is to concentrate proof effort on the narrow waist of the system: the
pure logic that decides whether a side effect may proceed and whether the
evidence describing that decision is structurally sound.

The resulting claim is intentionally modest but meaningful. The repository does
not prove the entire system correct. It does, however, provide a layered and
reproducible proof mesh over the logic whose failure would most directly break
the security boundary: fail-closed verdict computation, delegation attenuation,
approval safety, audit integrity, and normalization kernels. By pairing those
proofs with an explicit trust boundary and an explicit refinement story, the
artifact aims to make its strongest claims both useful and reviewable.

\bibliographystyle{plain}
\bibliography{references}

\end{document}
